\section{The quantum reduction}\label{app:quantumproof}
Let $H_M=\sum_E E\Pi_E$ and $[\rho,H_M]=[\tau,H_M]=0$. Each state is block diagonal in the subspaces $\Pi_E$. Choose separate unitaries $U_\rho$ and $U_\tau$, each block diagonal in those subspaces, that diagonalize the respective states in one fixed energy eigenbasis. Both commute with $H_M$. Their diagonal entries $p$ and $q$ have the same entropies and mean energies as $\rho$ and $\tau$.

Under the maximal gate rule these unitaries, their inverses and all classical energy-shell permutations are allowed. Given $U(\rho)=U(\tau)$ and $s(\rho)\le s(\tau)$, first apply $U_\rho$, then the constructor of Lemma~\ref{thm:closed}, and finally $U_\tau^\dagger$ to the released system. Unitary invariance preserves the trace-distance error. The classical constructor memory is stationary and returns exactly; the two endpoint rotations act only on the supplied system. This proves sufficiency in~\eqref{eq:qclosed}. Necessity was proved in Appendix~\ref{app:bounds}.

For adiabatic conversion, apply the same reasoning to the diagonalized states and the nine-level work construction of Appendix~\ref{app:calibration}. The endpoint rotations do not act on the work register, so its sharpness, dimension and span are unchanged. This proves~\eqref{eq:adiabaticcriterion} and the uniform bound~\eqref{eq:uniformmaxgate}. The argument applies to any finite composite as a single system. Its energy eigenspaces may contain entangled vectors.

For clarity, the two quantum gates in the running example are explicit. On the ordered basis $\ket{2,3},\ket{3,2}$ take
\[
 V_{\rm H}=\frac1{\sqrt2}\begin{pmatrix}1&1\\1&-1\end{pmatrix},\qquad
 Z_F=\begin{pmatrix}1&0\\0&-1\end{pmatrix},
\]
with the identity on the other 47 basis states. Both commute with the two-cell energy because the displayed block has energy five. The first maps $x$ to $\proj{\psi}$. The second, controlled by a uniform zero-energy bit, gives the joint state
\[
 \tfrac12\proj{\psi}\otimes\proj0+
 \tfrac12\proj{\psi_-}\otimes\proj1,
 \qquad \ket{\psi_-}=\frac{\ket{2,3}-\ket{3,2}}{\sqrt2}.
\]
Its system marginal is $\omega_{1/2}$ and its bit marginal is unchanged. Each cell's reduced state in $\proj{\psi}$ is
$\tfrac12(\proj2+\proj3)$, which is mixed. Since the joint state is pure, it cannot be a product and is entangled. No coherence between different total energies was used.

\section{Work, heat and information in the completion}\label{app:media}
The conversion criterion is now a proved consequence of the specified gates. We use it to check each operational clause of Appendix~\ref{app:clauses}, rather than using the clauses to stipulate conversion rules.

\subsection{Storage, preparation and locality}
A stationary state is fixed by free evolution generated by its Hamiltonian, so every admissible state can be stored indefinitely. At the gate level the identity also stores it. Generic preparation resources are sharp configurations, sharp work ladders and uniform degenerate bits. For a diagonal target $p$ on $d$ configurations, choose nonnegative integers $m_i$ summing to $2^N$ with $m_i/2^N\to p_i$. Divide the $2^N$ bit strings into sets of those sizes. Start the system in one sharp configuration. Conditional on the retained string, map it to the chosen output configuration and shift a work ladder by the compensating energy. A common rational spacing can accommodate every finite energy difference, with enough guard levels. This specified map is injective, because the string is retained, and preserves total energy on each branch. Extending the unused inputs within every shell makes it one energy-preserving permutation. It prepares the required marginal to arbitrary accuracy.

For a stationary quantum target, follow that preparation by its block-diagonal eigenbasis rotation. The bits, work store and other preparation outputs are consumed resources, not an omitted catalyst. A finite number of independent preparations uses separate finite supplies. A constructor memory state is prepared by the same procedure to any specified initial accuracy. The stationary-memory contraction argument in Appendix~\ref{app:clock} controls that initial error uniformly in the number of visits.

Locality can be stated without replacing a correlated state by its marginals. Purify any initial density operator on a mathematical reference and write the purified state as $V\ket0$. For each physical factor $A$, its Heisenberg descriptors are $V^\dagger(E_{ij}^{A}\otimes I)V$, where $E_{ij}^{A}$ are matrix units. Products of these descriptors recover all joint expectations in the fixed reference state $\ket0$. A physical unitary confined to $A$ commutes with every operator belonging to a distinct factor $B$, and therefore leaves $B$'s descriptors unchanged. The mathematical purification is a representation of the state, not an uncounted physical preparation resource. This verifies the stated locality clause also for entangled states.

\subsection{Information operations}
Perfect one-use discrimination of density operators requires orthogonal supports. To see necessity, a distinguishing effect must have expectation one in one state and zero in the other. Positivity then makes it the identity on the first support and zero on the second, which requires orthogonality. For approximate discrimination to arbitrary accuracy, trace-distance contraction gives the same conclusion: nonorthogonal supports have trace distance strictly less than one, leaving a positive error lower bound.

For stationary states with orthogonal supports, their support projectors $P_i$ commute with $H_M$. The block unitary
\[
 \sum_i P_i\otimes X_i+
 \Bigl(I-\sum_iP_i\Bigr)\otimes I
\]
copies the label to a blank degenerate record; $X_i$ is a permutation taking the blank label to $i$. The input state is undisturbed, the record is sharp for each promised input and the gate conserves energy. This is one finite circuit for the whole variable. Tensor products of such variables can be copied with product records. A finite pairwise-distinguishable family has pairwise orthogonal support projectors, so the same construction jointly distinguishes it, proving the union clause.

A permutation of the labels is possible with counted side effects. Copy the input label, prepare the finitely many prescribed output states on fresh identical registers using the generic resources just described, and use the label to swap the required register into the output position. Controlled swaps of identical Hamiltonians conserve energy; the input and unused preparation registers remain explicit side effects. Every branch is realized by the same finite circuit. Its finite accuracy can be chosen to meet all finitely many promises simultaneously. This proves the computation and copying requirements without asserting that nonorthogonal unknown quantum states can be cloned.

\subsection{Work tests and interoperability}
For an equally spaced sharp list, distinct isolated levels have unequal energies and hence cannot be interconverted by~\eqref{eq:qclosed}. Opposite adjacent steps on two replicas have the same total energy: the configurations $(j+1,k)$ and $(j,k+1)$ are exchanged by a transposition within one shell. Both directions are available. Thus adjacent steps are is-like.

For $1\le j\le L-1$, the two prescriptions in~\eqref{eq:workswap} have distinct sources, distinct targets and conserve energy separately. At the bottom use $(0,1)\mapsto(1,0)$ and $(1,1)\mapsto(0,2)$. Each partial injection extends to a permutation in each shell, by matching unused source and target configurations. This is one gate for both promised inputs. The top pair uses the first prescription with the replica initially at its lower member, and has the required downward guard. No endpoint outside $0,\ldots,L$ is used.

For two matched work lists of the same gap, the same prescriptions use the first substrate for the first coordinate and its matched replica levels on the second. Different additive energy offsets cancel. This verifies interoperability I(i) with a well-typed ancillary initialization. The product list has energies in an arithmetic progression with that gap. Choose one known product configuration at each distinct total energy. The identical shell-injection proof supplies its work tests, and configurations of repeated total energy can be exchanged directly. This proves I(ii) in the stated representative convention.

If a designated composite work pair factorizes, each changing factor is a known pure mechanical pair, and an unchanged factor is an identity. A pair with a nonzero rational energy difference exchanges that difference with a rational sharp work ladder; an equal-energy pair is connected by an energy-shell unitary. This is the factorization convention of Appendix~\ref{app:clauses}, including two-state factors. A fixed pure spectator neither changes these gates nor supplies unreported randomness.

\subsection{The four heat requirements and the zeroth law}
For a nonconstant finite Hamiltonian, $\gamma_\beta$ in~\eqref{eq:gibbs} is full rank. Let $m>0$ be its smallest eigenvalue. For any density operator $r$, $\gamma_\beta-mr\ge0$, since $\gamma_\beta\ge mI$ and $r\le I$. Consequently $\gamma_\beta=mr+(1-m)v$ for another state $v$, and $d(\gamma_\beta,r)\le1-m$. Tensoring an independently prepared apparatus leaves that distance unchanged; all subsequent processing contracts it. Two output records within $\eps$ of disjoint labels would require $1\le1-m+2\eps$, hence $\eps\ge m/2$. This proves the first heat requirement uniformly over all finite apparatus sizes and permitted side effects.

Differentiating the finite partition function gives $U=-\partial_\beta\log Z$ and $s(\gamma_\beta)=\beta U+\log Z$. The second derivative of $\log Z$ is the strictly positive energy variance. Thus~\eqref{eq:derivatives} holds and entropy strictly orders the family. Equations~\eqref{eq:adiabaticcriterion} and additivity give the required single direction for each distinct pair, both on one copy and on its identical-copy family. This checks the quasi-heat property and the second requirement.

The entropy range is the open interval from the logarithm of the ground-space dimension to $\log\dim M$. The average of any two member entropies lies between them and in that interval. Continuity and strict monotonicity give the unique $\beta_F$ in~\eqref{eq:equalization}. Equal total entropy and~\eqref{eq:adiabaticcriterion} permit both directions with bounded calibrated work. This is the third requirement.

For the fourth, suppose a heat state $\gamma_\beta$ initializes a fixed finite ancillary substrate in a work swap. In the upward branch the sharp work gain is $g>0$. Energy conservation and exact memory return require the ancillary final energy to tend to $U(\gamma_\beta)-g$. Entropy monotonicity on the combined work--ancilla input requires its final entropy, in the same limit, to be at least $s(\gamma_\beta)$: the work output approaches a pure state on a fixed finite register. But~\eqref{eq:kelvin} bounds that entropy by $s(\gamma_\beta)-\beta g$. This is a contradiction. The output ancillary state was not assumed thermal, and approximation cannot remove the positive gap on a fixed finite substrate. The same proof allows a variable work store with bounded dimension and span and prohibits positive limiting work extraction from a single Gibbs state. Direct closed heat-to-sharp-work conversion is also forbidden, by a strict entropy decrease.

Finally, $U(\gamma_\beta)$ decreases continuously from $\overline e$ to $e_{\min}$ as $\beta$ runs from zero to infinity. This proves the existence and uniqueness of heat counterparts on~\eqref{eq:interval}. At equal mean energy,~\eqref{eq:kelvin} makes the heat state's entropy maximal. Equation~\eqref{eq:qclosed} therefore proves X for every admissible state at that energy. Equal-energy work attributes have the same heat counterpart. No full-rank state can attain an extremal mean energy of a nonconstant Hamiltonian, because it puts positive probability on a nonextremal eigenspace. This checks the stated working-domain boundary.

\subsection{Why the rational work witness is calibrated}\label{app:fixedwork}
The departure from a single fixed work witness in heat equalization is necessary for the chosen rational-spectrum class. Fix one positive-temperature Gibbs state and vary a second on the same seven-level cell. The entropy-matching $\beta_F$ is continuous, as the inverse of a strictly monotone continuous function. The energy difference $W$ in~\eqref{eq:workengine} is therefore continuous. It vanishes when the two states agree. Substitution of their Gibbs logarithms, followed by~\eqref{eq:equalization}, gives
\[
 \DD(\gamma_{\beta_1}\otimes\gamma_{\beta_2}
       \Vert\gamma_{\beta_F}^{\otimes2})=\beta_F W.
\]
For unequal initial states the relative entropy is positive, so $W>0$. A continuous nonconstant $W$ takes all values in some interval, including irrational ones in the fixed energy unit.

A fixed pair of sharp work levels in a rational spectrum has a rational energy gap. Both directions of a closed compensating task require equality of the energy and entropy differences, by~\eqref{eq:qclosed}. Entropy equality already fixes $\beta_F$, while energy equality would require its $W$ to equal that rational gap. This fails for the irrational examples, even with arbitrary finite returned memory. Bounded calibration removes precisely that arithmetic obstruction. The argument concerns this rational-spectrum realization, not a prohibition on real energy gaps in other subsidiary theories.

\subsection{Heat interoperability and pair-task composition}
For pair tasks $A=(\rho,\tau)$ and $B=(\sigma,\upsilon)$, the closed is-like relation means that both $\rho\otimes\upsilon\plain\tau\otimes\sigma$ and its reverse are possible. Equation~\eqref{eq:qclosed} proves
\begin{equation}
 A\text{ is-like }B\quad\Longleftrightarrow\quad
 \begin{cases}
 U(\tau)-U(\rho)=U(\upsilon)-U(\sigma),\\
 s(\tau)-s(\rho)=s(\upsilon)-s(\sigma).
 \end{cases}
 \label{eq:islikeUH}
\end{equation}
This is a consequence of dynamics, not an assigned pair label.

Suppose a correspondence $f$ between two full heat families makes all corresponding pair tasks is-like. Fix one reference point. Equation~\eqref{eq:islikeUH} gives constants $a,b$ with $U_2(f(\beta))=U_1(\beta)+a$ and $s_2(f(\beta))=s_1(\beta)+b$, where the subscripts denote the family. The first equation makes $f=U_2^{-1}\circ(U_1+a)$ differentiable, since $U'_2$ never vanishes. Differentiating both equations gives
\[
 U'_2(f(\beta))f'(\beta)=U'_1(\beta),\qquad
 f(\beta)U'_2(f(\beta))f'(\beta)=\beta U'_1(\beta).
\]
Since $U'_1\ne0$, one has $f(\beta)=\beta$. The matched product is exactly the positive-temperature Gibbs family of $H_1+H_2$. The previous four checks apply to this finite composite, proving IX on the declared family domain.

Closed pair tasks compose because equality of mean energies and the entropy inequality in~\eqref{eq:qclosed} compose. Tensor products use additivity. The same reasoning applies to the adiabatic relation in~\eqref{eq:adiabaticcriterion}. The existence proof constructs a suitable finite memory for the resulting pair task. It does not assert that the product of two arbitrary catalysts returns after their memories have become correlated. Comparison follows from the order of real entropies, and reflexivity follows from the identity.

Every equal-energy state reaches the Gibbs maximizer at that energy, or the uniform state in an extremal eigenspace, by~\eqref{eq:qclosed}. Hence any function conserved by every closed pair task is constant on a fixed-energy set. There is no second independent state charge. A closed energy mismatch forbids both directions; at equal energy the only obstruction is decreasing entropy. Different work levels instantiate the first obstruction, while $\delta_1\plain u$ with its forbidden reverse instantiates the single irreversibility condition. This completes the work, heat, information, conservation, preparation, existence and pair-composition checks for Theorem~\ref{thm:quantum}.

\section{The different product-return contract}\label{app:product}
Suppose~\eqref{eq:return} is strengthened to $R=R_M\otimes c$ exactly, rather than just $R_C=c$. Unitary entropy balance then gives $s(R_M)=s(\rho)$. For an adiabatic witness with exact independence across $MB:C$, the same calculation gives $s(R_{MB})=s(\rho)$. Uniform entropy continuity under bounded work dimension forces $s(\tau)=s(\rho)$ in the accuracy limit. Thus entropy-changing pair tasks are forbidden in both directions. On the seven-level cell, $\delta_1$ and $u$ from~\eqref{eq:smallclock} have the same energy but entropies zero and $\log3$. They are incomparable under exact product return. The 1029-configuration circuit in Appendix~\ref{app:clock} permits the increasing direction with marginal return and explicitly displays the correlations responsible. This is the product-return obstruction already emphasized in the correlated-catalysis literature~\cite{Boes,Wilming2021}.

\section{Coherence and comparison}\label{app:coherenceproof}
We state the prior no-broadcasting result in the form needed here and include its finite structural reduction. The non-disturbance structure itself is the Koashi--Imoto theorem~\cite{KoashiImoto}, not a new result of this paper. This proof is the time-translation specialization of Lostaglio--M\"uller's argument~\cite{LostaglioMueller}; the independent theorem of Marvian--Spekkens gives the same obstruction~\cite{MarvianSpekkens}.

\begin{lemma}[Finite-reference no broadcasting]\label{lem:nobroadcast}
Let $S$ start in $\sigma$ commuting with $H_S$, and let a finite reference $C$ start in an arbitrary state $c$. If an energy-preserving unitary maps $\sigma\otimes c$ to $R$ with $R_C=c$, then $[R_S,H_S]=0$. Correlations between $S$ and $C$ are allowed.
\end{lemma}
\begin{proof}
For real $t$ put $c_t=e^{-itH_C}ce^{itH_C}$. Covariance and stationarity of $\sigma$ imply that the channel $\Lambda(a)=V(\sigma\otimes a)V^\dagger$ returns the reference marginal of every $c_t$ unchanged. Its system output is $e^{-itH_S}R_Se^{itH_S}$. Restrict the reference space to the span of the supports of the $c_t$; all states and outputs stay in this finite space.

The finite non-disturbance structure~\cite{KoashiImoto} says that a family left invariant by the marginal channel has a fixed decomposition
\[
 C=\bigoplus_j L_j\otimes R_j,\qquad
 c_t=\bigoplus_j p_j(t)\,a_j(t)\otimes b_j.
\]
Here $b_j$ is independent of $t$. In a dilation of a channel that leaves this family unchanged, the environment can depend on the block label $j$ and its fixed redundant state $b_j$, but not on the variable state $a_j(t)$. In particular, the $S$ output has the form $\sum_j p_j(t)\zeta_j$ with $\zeta_j$ independent of $t$.

Each displayed block of $c_t$ is a continuously varying Hermitian matrix. Its eigenvalues belong to the fixed finite spectrum of $c$, because $c_t$ is block diagonal and unitarily equivalent to $c$. Every ordered block eigenvalue is therefore a continuous function from the connected real line to a finite set, and is constant. Its block trace $p_j(t)$ is constant as well. The system output is consequently independent of $t$. Thus $e^{-itH_S}R_Se^{itH_S}=R_S$ for every $t$, which is equivalent to the claimed commutation.
\end{proof}

Apply the lemma with $S=MB_\eps$. The input $g\otimes w_i$ is stationary because both factors are sharp energy states. The exact whole memory, including a coherent reference if supplied, is $C$. Each output $R_{MB_\eps}$ is stationary, and partial trace makes $R_M$ stationary. The set of stationary states is closed, since commutation with a fixed finite matrix is continuous. No accuracy sequence can approach a nonstationary target.

For the qubit target~\eqref{eq:coherentqubit}, a stationary density operator is $\operatorname{diag}(a,1-a)$. Its difference from $\rho_\chi$ has eigenvalues
$\pm\sqrt{(a-1/2)^2+\chi^2/4}$, so the minimum trace distance is $\chi/2$. The eigenvalues $(1+\chi)/2,(1-\chi)/2$ of $\rho_\chi$ are both positive, hence $s(\rho_\chi)>0$. Monotonicity forbids $\rho_\chi\ad g$. This proves the stated explicit incomparability.

Finally, suppose a convex admissible domain contains a nonstationary $\rho$, the maximally mixed state and a pure energy state $g$. For any $0<\alpha<1$, the state $\sigma=(1-\alpha)\rho+\alpha I/\dim M$ remains nonstationary and has full rank. Since a nonstationary state requires dimension at least two, $s(\sigma)>0$. The lemma forbids $g\ad\sigma$ and entropy monotonicity forbids $\sigma\ad g$. This contradicts comparison and proves Theorem~\ref{thm:domain}. The stationary completion is therefore maximal within the stated convex class of domains.
